\section{Executable Evidence}
\label{sec:evidence}

\subsection{Method}
\label{sec:method}

\paragraph{Kernel and simulator}
The kernel (\file{src/mandate/}) implements O1 to O3, D1 and persistence as functions over
explicit state, with a canonical message encoding. Signatures have two modes: a symbolic
mode, and Ed25519~\cite{rfc8032} with fixture keys that are only for tests. A file store
makes state durable before a signature is released. The simulator (\file{sim/}) runs a
\emph{scenario}, which is a list of timed events: signatures, deliveries, commands,
withdrawals, compromises, crashes and restarts. Events with equal times run in list
order. Each party has a constant clock offset of at most $\eps$. The output is a trace in
canonical JSON, one event per line, with a SHA-256 digest. For every scenario the two
signature modes must give byte-identical traces, and they do.

\paragraph{Independent checker}
The checker (\file{experiments/checker.py}) reads only the trace and imports nothing
from the kernel or the simulator. For each acceptance it checks the evidence again: a
grant and a renewal, each signed by $q$ distinct configured officers whose signing events
appear earlier in the trace, and delivered before the acceptance. It reports grant
uniqueness, epoch order, the pairwise separation of the 11 November draft, replay,
invalid evidence, the withdrawal-to-expiry bound, clock bounds, trace order, and two
holders accepted at one clock reading. It also computes metrics such as $t_W$ and the
last acceptance of a compromised holder. The checker is independent in code but not in
authorship: the same assistant wrote both, so a shared misreading of the specification
would pass.

\paragraph{Expectations and manifest}
Each scenario states its expected violations exactly, together with expected metric
values. A run fails if any expectation fails. The manifest records the Git commit and
whether the tree was dirty, a SHA-256 over the source files, the Python and library
versions, and the digest of every configuration and trace. There is no randomness, and
all times are synthetic integers.

\subsection{Directed Scenarios}
\label{sec:scenarios}

\Cref{tab:scenarios} lists the scenarios that exercise the rules. The positive scenarios
use $n = \SeedN$, $f = \SeedF$, $q = \SeedQ$, $H = \SeedH$, $g = \SeedG$ and
$\eps = \SeedEps$, so $\gamma = \SeedGamma$. All \SeedScenarioCount{} scenarios meet their
expectations (\SeedPassCount{} of \SeedScenarioCount{}).

\begin{table*}[t]
\caption{Directed scenarios of seed-v0 and of the 11 November 2026 rules, with the
negative controls. ``Acc.'' is the number of accepted commands. ``Violations found'' are
the checker's codes. ``Met'' means that every expectation of the scenario holds.
Generated from the traces.}
\label{tab:scenarios}
\centering
\footnotesize
% Generated by experiments/report.py. Do not edit by hand.
\begin{tabular}{@{}l>{\raggedright\arraybackslash}p{0.33\linewidth}rll@{}}
\toprule
Scenario & What it exercises & Acc. & Violations found & Met \\
\midrule
\texttt{V0-RENEW-EXPIRY} & renewal, withdrawal, expiry at two verifiers & 5 & none & yes \\
\texttt{V0-HANDOVER} & prepared handover; separation is tight ($\gamma=1$) & 6 & none & yes \\
\texttt{V0-SELECTIVE-DELIVERY} & delayed and selective delivery to two verifiers & 6 & none & yes \\
\texttt{V0-REMOVAL} & hostile holder, withdrawal race, successor & 6 & none & yes \\
\texttt{V0-ISOLATION} & isolated honest holder; old command delivered late & 5 & none & yes \\
\texttt{V0-REPLAY-CLOCK} & replay, boundaries, equal-time order, forged evidence & 7 & none & yes \\
\texttt{V0-SPLIT-RECOVERY} & split epoch-2 vote, stall, epoch-3 recovery & 3 & none & yes \\
\texttt{ORIG-SPLIT-STALL} & draft rules: split vote stalls service & 2 & none & yes \\
\texttt{ORIG-SPLIT-RECOVER-E3} & draft rules: same stall, epoch-3 recovery & 3 & none & yes \\
\texttt{ORIG-DEPUTY-SAME-READING} & draft rules, $F>0$: two issuers at one reading & 3 & SAME-READING, SEPARATION & yes \\
\midrule
6 \texttt{NC-*} & broken rule variants (one guard removed each) & --- & expected code & 6/6 \\
6 \texttt{*-CONTROL} & the same events with the \texttt{seed-v0} rules & --- & none & 6/6 \\
\bottomrule
\end{tabular}

\end{table*}

\paragraph{What the scenarios show}
In \texttt{V0-REMOVAL} (\cref{fig:removal}), a renewal that a withdrawing agent signed just
before $t_W$ keeps the hostile holder accepted until $t_W + \RemovalAfterTW$, within the
candidate bound $t_W + H + 2\eps = \RemovalBound$. In \texttt{V0-ISOLATION}, renewal
requests start every $\Delta = \IsolationDelta$ and complete within
$\delta = \IsolationCollect$, and the outage starts at \IsolationOutage{}. A fresh command
at the end of the guaranteed window (\IsolationWindowEnd) is accepted. So is a later one,
\IsolationFreshSpan{} after the outage starts. A command created before the outage and
delivered late is accepted as well, but it does not count toward isolation reach.
\texttt{V0-HANDOVER} shows that separation is tight: the gap is \HandoverGap{}, which is
$\gamma$. \texttt{V0-REPLAY-CLOCK} covers replay, both window boundaries, the order of
equal-time events, and forged and cross-scope evidence.

\paragraph{The draft's rules}
Two fixtures pin findings about the draft of 11 November 2026. In
\texttt{ORIG-SPLIT-STALL}, a split vote leaves epoch 2 without a valid grant, and no
command is accepted after the last renewal expires. No safety property is violated. A
coordinated epoch-3 grant restores service (\texttt{ORIG-SPLIT-RECOVER-E3}), and seed-v0
behaves the same (\texttt{V0-SPLIT-RECOVERY}). In \texttt{ORIG-DEPUTY-SAME-READING}, the
deputy rule ($F > 0$) gives two issuers at one clock reading. The verifier accepts the
deputy, then receives a renewal that the adversary held back, then accepts the holder.
Part~1 of the draft's deputy proposition (``at one clock time, one device accepts orders
of not more than one officer'') is therefore false \claim{C-DEP}. The draft's proof checked
both rules in one verifier state, but the state changed between the two acceptances.

\paragraph{Negative controls}
Each of six broken variants removes one guard: the O2 start guard, the O1 epoch guard,
D1's end check, D1's replay check, distinct-signer counting, or persistence before
release. The checker reports the intended violation for \SeedNegativeDetected{} of
\SeedNegativeCount{}. The control twins run the same events under seed-v0, and
\SeedControlClean{} of \SeedNegativeCount{} are clean \claim{C-EV-NC}. This shows that the checker can see
each fault in its scenario, not that it finds every fault of that kind.

\subsection{Comparison Mechanisms}
\label{sec:comparisons}

We defined three mechanism abstractions. None of them reproduces a published system.
\begin{description}
  \item[\texttt{tc}] Threshold-issued expiring credentials. A quorum signs
  $\mathsf{Lease}\langle h, s, e\rangle$ with $e \le C_u + H$. There is no exclusive
  succession rule.
  \item[\texttt{xl}] Exclusive quorum leases, our Byzantine adaptation of Nerio's grantor
  rule. Agent $u$ keeps its current holder $A_u$, the start $S_u$ of its first lease for
  $A_u$, and its largest signed end $F_u$. It signs $\mathsf{Lease}\langle h, s, e\rangle$ only if it consents to
  $h$, $s \le e \le C_u + H$, and either $h = A_u$ and $s \ge S_u$, or $h \ne A_u$ and
  $s \ge F_u + g$.
  \item[\texttt{qo-w}] A quorum signs every command, with a validity window of length $W$.
\end{description}
All mechanisms run the same command stream in each comparison. \Cref{tab:comparison}
gives the results.

\begin{table}[t]
\caption{Comparison runs (one directed run per cell). Handover: time from the old holder's
last acceptance to the new holder's first, or ``overlap'' if the checker found
overlapping holders. After $t_W$: the compromised holder's last acceptance minus $t_W$,
with the candidate bound in brackets. Isolated: how long after the outage started the
latest accepted fresh command was created. Split: whether the old holder keeps service
after a split vote for the same holder, and for two holders.}
\label{tab:comparison}
\centering
\footnotesize
\setlength{\tabcolsep}{2.5pt}
% Generated by experiments/report.py. Do not edit by hand.
\begin{tabular}{@{}lllll@{}}
\toprule
Mechanism & Handover & After $t_W$ & Isolated & Split: same / two \\
\midrule
\texttt{seed-v0} & gap 16 & 98 ($\le 102$) & 92 & stalls / stalls \\
\texttt{xl} & gap 16 & 98 ($\le 102$) & 92 & continues / stalls \\
\texttt{tc}, no wait & overlap & 98 ($\le 102$) & 92 & not run \\
\texttt{tc}, wait & gap 16 & as above & as above & not run \\
\texttt{qo-w}, $W=20$ & not run & 19 ($\le 22$) & none & not run \\
\bottomrule
\end{tabular}

\end{table}

\paragraph{Reading the comparison}
Without a waiting rule, \texttt{tc} overlaps two holders in the handover. If officers
wait for the old credential to expire, it does not, but then exclusivity rests on their
discipline and not on a signer rule. \texttt{xl} gives the same values as seed-v0 in the
handover, removal and isolation runs. In a split vote for the same holder, \texttt{xl}
continues and seed-v0 stalls. The cause is O2, which advances $E_u$ when an agent votes.
\texttt{qo-w} bounds in-flight hostile commands by $W + 2\eps$ after $t_W$, but accepts no
fresh command in isolation \claim{C-CMP}. Equal values come partly from the shared
schedules, and one run per cell measures nothing about cost.

\subsection{Bounded Model Checking}
\label{sec:tla}

The TLA+ model~\cite{tlaplus} (\file{spec/tla/}) covers O1 to O3 and D1 with
\TlaHonest{} honest signers and \TlaByz{} Byzantine signer ($q = \TlaQ$), \TlaHolders{}
holders, epochs up to \TlaMaxEpoch, time up to \TlaMaxTime{}, $H = \TlaH$ and $g = \TlaG$.
It assumes exact clocks, Byzantine signatures on every message, and a verifier that
receives every valid certificate at once. TLC~\TlaVersion{} found no violation of grant
uniqueness, epoch order, one holder per epoch or withdrawal-to-expiry \claim{C-TLA}. It explored
\TlaDistinct{} distinct states to depth \TlaDepth. Three further configurations must fail,
and they do. A handover is reachable (trace of \TlaHandoverTrace{} states), and so is the
split-vote stall (\TlaStallTrace{} states). Without the O2 start guard, epoch order is
violated (\TlaBrokenTrace{} states). A run with larger bounds was stopped before it finished. The
model's other abstractions are listed with it.

\subsection{What the Evidence Does Not Show}

Directed runs show that behaviours exist. They do not explore schedules: no randomized
or exhaustive search over the simulator was run. The bounded model has $\eps = 0$ and tiny
bounds. No run measures time, throughput or signature cost on real systems. The
parameters are synthetic, and nobody has tied them to an operational requirement.
